\section{K-031: the chain is the only cyclically stable fixed locus}

This strengthens K-026, K-028, and K-030 by removing the
stabilizer-dimension hypothesis.  It leaves the separate collision
exception identified in K-027 and explicitly realized in K-029.

\begin{theorem}[K-031a]
Let $n\ge4$, $M=2^{n-1}$, and $\tau=(123)$.  Every irreducible
locally closed subvariety $W\subset (X_n^{\mu_M})_{\rm red}$
preserved by $\tau$ is contained in the chain curve $Y_n$.
Consequently $Y_n$ is the only $\tau$-stable irreducible component
of this fixed locus.  Every other component has $S_3$ orbit size
three or six, without any assumption on its generic stabilizer.
Moreover,
\[
       \bigl((X_n^{\mu_{M3^a}})_{\rm red}\bigr)^{\tau}
                             =\{0\}\qquad(a\ge1)
                                                        \tag{K-031a}
\]
as sets of closed points.
\end{theorem}

\begin{proof}
There is nothing to prove if $W=\{0\}$.  Choose a positive-degree
homogeneous special-linear invariant nonzero on a dense subset of
$W$.  Replace it by the product of its three cyclic translates;
call the result $f$.  Then $f$ is homogeneous, $\tau$-invariant,
and still nonzero on a dense subset.  It is a semi-invariant for
$\overline G=GL_n^3/S$ of some positive character $\chi^a$.
Thus the principal open $V_f\subset V_n$ is smooth, affine, and
stable under both $\overline G$ and $\tau$.

\emph{Closed full orbits on this chart.}
If $T\in V_f$ is a closed special-linear orbit representative,
its full $\overline G$ orbit is closed in $V_f$.  Here is a
direct proof, including positive-dimensional stabilizers.  Put
$y=q(T)$, and let $d$ be the dimension of its effective
special-linear orbit.  The nonzero grading orbit of $y$ has
only the origin in its affine boundary and is therefore closed
on the quotient chart $f\ne0$.  Every stabilizer on $V_f$ has
finite $\chi$ image, since it fixes the nonzero value of $f$.
Hence the full orbit of $T$ has dimension $d+1$.

If a point $T'$ were in its boundary inside $V_f$, its quotient
would lie on the same grading orbit.  The unique closed
special-linear orbit in that quotient fibre is then a scalar
translate of the orbit of $T$.  Since $T'$ is outside the full
orbit, its special-linear orbit is nonclosed and has dimension
strictly greater than $d$.  Its full orbit consequently has
dimension strictly greater than $d+1$.  This is impossible for
an orbit in the boundary of a $(d+1)$-dimensional orbit.
The full orbit is therefore closed, as asserted.

Apply Luna's finite model stratification to the smooth affine
$\overline G$-variety $V_f$.  Section III.2, Corollaries 3--4
of Luna (1973), printed p.~100, gives finitely many locally
closed model strata, hence finitely many full-stabilizer
conjugacy types for closed orbits.  The quotient map by
$\overline G$ factors through the special-linear quotient
chart, so these types pull back to a constructible partition
of $W\cap\{f\ne0\}$.  After shrinking to a dense open subset,
the full stabilizer modulo $S$ has one ambient conjugacy type
$H$.  The shrinking can be made $\tau$-stable.  It follows that
$\tau(H)$ is conjugate to $H$ in $\overline G$.

\emph{A finite subgroup carrying both characters.}
Let $L=N_{\overline G\rtimes\langle\tau\rangle}(H)$.
Its projection to $\langle\tau\rangle$ is surjective.  It
contains a finite three-primary element $r$ projecting to
$\tau$.  To see this without a reductivity assertion for the
normalizer, choose any element above $\tau$ and take its
semisimple part.  Its unipotent part maps trivially to the
finite quotient.  The Zariski closure of the cyclic group
generated by the semisimple part is diagonalizable.  Its
finitely generated character group splits as a free group
plus torsion, so this diagonalizable group splits as a torus
times a finite group.  The finite three-primary factor
surjects onto $\langle\tau\rangle$, giving the desired $r$.
This is the same arbitrary-order finite lifting argument used
in K-009; no order-three lift is assumed.

Set $J=H\langle r\rangle$.  This is a linear algebraic group
which is a finite union of cosets of $H$, and $J^\circ=H^\circ$.
Its determinant character has finite image: $\chi(H)$ is finite
on $V_f$, and $r$ has finite order.  In particular
$\chi(H^\circ)=1$.

Use the classical finite-subgroup theorem of Borel--Serre
(1964), Lemma 5.11, printed p.~152: a complex linear algebraic
group contains a finite subgroup meeting every connected
component.  A primary modern statement sufficient
here is Lucchini Arteche (2017), Theorem 3.2, applied to
$1\to J^\circ\to J\to\pi_0(J)\to1$.  It gives a finite
subgroup $F\subset J$ surjecting onto $\pi_0(J)$.
Then $N=F\cap H$ is a normal subgroup of $F$ and surjects onto
$\pi_0(H)$, since $H$ is a union of components of $J$.
Consequently
\[
                       \chi(N)=\chi(H).
                                                        \tag{K-031b}
\]
Also $F$ maps onto $\langle\tau\rangle$.  A Sylow
three-subgroup of $F$ has the same property, so it contains
an element $b$ projecting to $\tau$.

Fix a primitive $M$th root $\xi$.  The fixed-point condition
gives $\xi\in\chi(H)=\chi(N)$.  Conjugation by $b$ induces
an automorphism of three-power order on the finite group $N$
and preserves $\chi$.  Apply K-030a to obtain a two-primary
$h\in N$ with $\chi(h)=\xi$ whose $N$-conjugacy class is
preserved by that automorphism.  Since $b$ has component $\tau$
in the semidirect product, the ambient $\overline G$-conjugacy
class of $h$ is preserved by $\tau$.

\emph{Support classification.}
The remaining argument no longer uses the dimension of $H$.
Lift the finite two-primary element $h$ to a finite two-primary
triple in $GL_n^3$ by taking roots in $S$.  Cyclic stability
of its ambient class makes its three spectra scalar translates.
As three is invertible modulo its two-power order, scalar shifts
of total zero make them equal.  After independent conjugation
the fixing element is $(A,A,A)$, with $A$ of two-power order
and $(\det A)^3=\xi$.  Thus $\det A$ has exact order $M$.

For every representative in the generic stabilizer-type locus,
a conjugate of this element fixes it.  Product-torus
semistability, cubic folding, the complete supplied theorem
P-F30, and the unique positive chain margins give precisely
the same conclusion as in K-026: that tensor is product-linearly
equivalent to the chain.  A dense subset of $W$ is contained
in $Y_n$, so all of $W$ is contained in the closed curve $Y_n$.

The statement about component orbits follows as in K-026.
For the final assertion, a nonzero $\tau$-fixed point in
$X_n^{\mu_{M3^a}}$ would, by the theorem applied to its one-point
subvariety, lie on $Y_n$.  But every nonzero point of $Y_n$ has
grading stabilizer exactly $\mu_M$, by P-F60.  This contradiction
proves (K-031a).
\end{proof}

\begin{corollary}[K-031c: collision closures are the sole exceptional support]
In the integral Laurent reduction K-027, one can take the
exceptional class to be supported on
\[
 Z_{\rm coll}=\{0\}\ \cup\!
 \bigcup_{\substack{Y_i\ne Y_n\ \text{an irreducible component}\\
                  Y_i\cap X_n^{\mu_{3M}}\not\subseteq\{0\}}}Y_i.
                                                        \tag{K-031c}
\]
There is no additional stabilizer exception.  If the exact
primitive-$M$ pole cancels, this collision-supported virtual
class must have the nonzero principal-part congruence specified
in K-027.
\end{corollary}

\begin{proof}
Repeat K-028b's finite dimension induction in the
$\mathbb G_m\times\langle\tau\rangle$-equivariant coherent
Grothendieck group, using K-031a at every residual irreducible
stratum.  Outside the chain and $Z_{\rm coll}$, every such
stratum occurs in a triple.  Its coherent extensions are
supported in closures contained in original components with
no nonzero $\mu_{3M}$ fixed point.  Thus every Laurent
coefficient of their Hilbert terms is integral at three by
K-027's supported-module no-collision lemma.  All induced
triples contribute multiples of three, leaving the chain
unit and a class supported on $Z_{\rm coll}$.  This is exactly
the proof of the stated Laurent reduction.
\end{proof}

\paragraph{Primary source and proof scope.}
Lucchini Arteche, G. (2017). Extensions of algebraic groups with
finite quotient and nonabelian 2-cohomology. \emph{Journal of
Algebra, 492}, 102--129. doi:10.1016/j.jalgebra.2017.08.026.
Theorem 3.2 is on printed p.~18 of arXiv:1503.06582v3; its
complete proof, Steps 1--4 on pp.~19--23, and the preceding
Lemma 3.3 and proof were read for this application.  The original
Borel--Serre Lemma 5.11 and its complete proof on printed p.~152
were also read directly: Borel, A., and Serre, J.-P. (1964),
\emph{Th\'eor\`emes de finitude en cohomologie galoisienne},
\emph{Commentarii Mathematici Helvetici, 39}, 111--164.
No novelty is claimed for the finite-subgroup theorem.  Only the existence
part of Theorem 3.2 is used here, over the perfect field
$\mathbb C$, where all groups involved are smooth.

K-029 proves that $Z_{\rm coll}$ contains actual large amplified
components in infinitely many dimensions.  K-031 does not
calculate their Hilbert contribution, and the exact primitive-$M$
tensor pole remains a separate unresolved question.
